
\emph{est de type fini.}

En effet, on peut supposer \'evidemment \(S\) r\'eduit et irr\'eductible d'apr\`es 3.4 (i). Alors il existe un ouvert non-vide \(U\) de \(S\) tel que \(\underline{\Pic}_{Y/S}\) et \(\underline{\Pic}_{X/S}\) soient repr\'esentables et que la restriction de \(f^*\) soit quasi-affine (XII 1.1), donc de type fini par 3.2 (i). Munissant \(T=S-U\) de la structure r\'eduite induite, on peut supposer que le morphisme \(\underline{\Pic}_{Y_T/T}\longrightarrow\underline{\Pic}_{X_T/T}\) induit par \(f^*\) est de type fini par r\'ecurrence noeth\'erienne\,; d'o\`u l'assertion, gr\^ace \`a 3.4 (i) appliqu\'e \`a \(U\amalg T\to S\).

\textbf{Th\'eor\`eme 3.6.} \emph{Soient \(S\) un sch\'ema noeth\'erien, \(X\) un \(S\)-sch\'ema propre et \(\underline{\Pic}_{X/S}\) le foncteur de Picard. Alors pour tout entier \(n\ne0\), l'homomorphisme de puissance \(n\)-i\`eme}
\[
\Phi_n:\underline{\Pic}_{X/S}\longrightarrow\underline{\Pic}_{X/S}
\]
\emph{est un morphisme de type fini.}

En effet, soit \(f:X'\to X\) une pr\'esentation de Chow~: \(X'\) est projectif sur \(S\) et \(f\) est surjectif. Si l'on admet le th\'eor\`eme pour \(X'\), alors en vertu de 3.5 et 3.4 (ii), le th\'eor\`eme pour \(X\) en r\'esulte. On peut donc supposer \(X/S\) projective.

Il existe un ouvert non-vide \(U\) de \(S\), un sch\'ema int\`egre \(S'\) et un morphisme fini et surjectif \(S'\to U\), tels que toute composante irr\'eductible \(X'_i\) de \(X_{S'}\), munie de la structure r\'eduite induite, soit int\`egre sur \(S'\) (XII 1.3). Munissant \(T=S-U\) de la structure r\'eduite induite, on peut admettre le th\'eor\`eme pour \(X_T/T\) par r\'ecurrence noeth\'erienne. Si l'on l'admet pour tous les \(X'_i/S'\), il en r\'esulte pour \(\coprod X'_i/S'\) par le
